% --------------------------------------------------------------------------------------- %
%                   UOC - CSD - HY119 - Spring 2025 - Ass2 - csd5748 - :)                 %
% --------------------------------------------------------------------------------------- %

\documentclass[fleqn]{article}
\usepackage{amsmath, amssymb}
\setlength{\mathindent}{0pt}
\usepackage[utf8]{inputenc} 
\usepackage[english, greek]{babel}
\usepackage{subcaption}
\usepackage{hyperref}
\usepackage{xcolor}
\hypersetup
{
    hidelinks,  
    colorlinks=true,
    linkcolor=blue,
    urlcolor=blue 
}
\captionsetup[subfigure]{labelformat=empty}
\usepackage{amsmath}
\usepackage{amsfonts}
\usepackage{amssymb}
\usepackage{array}
\usepackage{graphicx}
\usepackage[a4paper, left=1.1cm, right=1.1cm, top=1.3cm, bottom=2cm]{geometry}
\usepackage[absolute, overlay]{textpos}
\usepackage{placeins}
\usepackage{listings}
\usepackage{xcolor}
\setlength{\jot}{0pt}
\lstdefinestyle{mypython}
{
    language=Python,
    basicstyle=\ttfamily\small,
    keywordstyle=\color{blue},
    stringstyle=\color{red},
    commentstyle=\color{green!50!black},
    numbers=left,
    numberstyle=\tiny\color{gray},
    stepnumber=1,
    frame=single,
    backgroundcolor=\color{gray!10},
    breaklines=true
}
\usepackage{cancel}
\setlength{\mathsurround}{0.1cm}
\setlength{\jot}{0.1cm}
\newcommand{\skato}{\hspace{0.1cm}}
\newcommand{\skata}{\hspace{0.321cm}}
\begin{document}
\selectlanguage{greek}



% ---------------------------------------------------------------------------------------



\begin{center}
\boxed{\textlatin{csd5748}}
\end{center}
\vspace{-0.735cm}
\begin{center}
\setlength{\fboxsep}{20pt}
\noindent{\boxed{{{\textbf{\Large{ΗΥ-119 - Γραμμική Άλγεβρα -  Άσκηση 2.}}}}}}
\end{center}



% ---------------------------------------------------------------------------------------



\vspace{1.23456789cm}
\begin{center}
    {{\textbf{·  \underline{1\textsuperscript{η} Άσκηση:}}}}
\end{center}
\vspace{0.54321cm}

\[
\textbf{Έχουμε:} \hspace{0.54321cm} 
A \skato = \skato
\begin{bmatrix}
& 1 & k & (4k+1) &  \\
& 0 & (1-k) & -k &  \\
& 0 & 0 & -k &
\end{bmatrix}
\skata \textbf{kai} \skata
b \skato = \skato
\begin{bmatrix}
& 4k+1 &  \\
& -k-1 &  \\
& -k &
\end{bmatrix}
\]


\vspace{0.75cm}
(α) \textbf{Για να είναι ο χώρος στηλών του πίνακα Α μια ευθεία στον \textlatin{R}3 θα πρέπει όλες οι στήλες του πίνακα να είναι γραμμικά εξαρτημένες μεταξύ τους και να ανήκουν σε έναν μονοδιάστατο χώρο, ή} \skato \textlatin{rank(A)} = 1. \textbf{Όμως,} \skato \textlatin{rank} = οδηγοί = 3. Αφού οδηγοί = 1, 1-\textlatin{k}, -\textlatin{k} θα πρέπει:
1=0 (αδύνατο), 1-\textlatin{k}=0, -\textlatin{k}=0 \skato \Rightarrow
\textlatin{k}=1 \textbf{ και } \textlatin{k}=0 \textbf{ (άτοπο)}. \textbf{ Τελικά, δεν υπάρχει τέτοιο } \textlatin{k!}


\vspace{1cm}
(β) Για να βρούμε ένα \textlatin{k} ανήκει \textlatin{R} για τον πίνακα Α ώστε ο μηδενόχωρος \textlatin{N(A)} να έχει διάσταση \textlatin{dim(N(A))}=1 πρέπει \textlatin{rank(A)}=2 αφού ξέρουμε: \textlatin{dim(N(A))=n-r=}στήλες πίνακα - αριθμός οδηγών. Άρα: \textlatin{dim(N(A))}=3-\textlatin{rank(A)} \Rightarrow 1=3-\textlatin{rank(A)} \Rightarrow \textlatin{rank(A)}=2. \textbf{ Έτσι, έχουμε πως οδηγοί}=1,1-κ,-κ. \textbf{ Οπότε, για } κ=0 \textbf{ και } κ=1 \textbf{ ισχύει: οδηγοί}=1,1 \textbf{ ή } 1,-1 \Rightarrow \textbf{ οδηγοί}=2 \Rightarrow rank(A) = 2. \textbf{ Και:}

\[
Ax=b: \skata
\begin{bmatrix}
& 1 & k & 4k+1 & | & 4k+1 & \\
& 0 & 1-k & -k & | & -k-1 & \\
& 0 & 0 & -k & | & -k & 
\end{bmatrix}
\Rightarrow
\left\{
\begin{array}{l}
k=0: 
\begin{bmatrix}
& 1 & 0 & 1 & | & 1 & \\
& 0 & 1 & 0 & | & -1 & \\
& 0 & 0 & 0 & | & 0 & 
\end{bmatrix}
\Rightarrow
\left\{
\begin{array}{l}
x_1 + x_3 = 1 \\
x_2 = -1 \\
0 = 0 \textbf{ (ισχύει)}
\end{array}
\right.
\skata *
\\
k=1:
\begin{bmatrix}
& 1 & 1 & 5 & | & 5 & \\
& 0 & 0 & -1 & | & -2 & \\
& 0 & 0 & -1 & | & -1 & 
\end{bmatrix}
\Rightarrow
\left\{
\begin{array}{l}
x_1 + x_2 + 5x_3 = 5 \\
-x_3 = -2 \\
-x_3 = -1 \textbf{ (αδύνατο)}
\end{array}
\right.
\end{array}
\right.
\]

\[
* \skata \skato
x_{\textbf{ειδ.}} = 
\begin{bmatrix}
1 \\
-1 \\
0
\end{bmatrix}
\skata \textbf{kai} \skata
x_{\textbf{γεν.}} = x_{\textbf{ειδ.}} + c \cdot 
\begin{bmatrix}
-F \\ I
\end{bmatrix}
=
\begin{bmatrix}
1 \\
-1 \\
0
\end{bmatrix}
+
C \cdot
\begin{bmatrix}
- \begin{bmatrix}
1 \\ 0
\end{bmatrix} \\ 1
\end{bmatrix}
=
\begin{bmatrix}
1 \\
-1 \\
0
\end{bmatrix}
+
C \cdot
\begin{bmatrix}
-1 \\
0 \\
1
\end{bmatrix}
=
\begin{bmatrix}
1-c \\
-1 \\
c
\end{bmatrix}
=
x_{\textbf{γεν.}}
\]

\vspace{1cm}
(γ) \skato Για \textlatin{k}=1 έχουμε:
A =
\begin{bmatrix}
& \boxed{1} & 1 & 5 &  \\
& 0 & 0 & -1 & \\
& 0 & 0 & \boxed{-1} & 
\end{bmatrix}
\Rightarrow
rank(A) = 2.
\textbf{ Οπότε: }  

\[
- \skato dim(C(A^t))=rank(A)=2
\]
\[
    - \skato dim(N(A)) = columns(A) - rank(A) = 3-2 = 1
\]
\[
- \skato dim(C(A))=rank(A)=2
\]
\[
    - \skato dim(N(A^t)) = rows(A^t) - rank(A) = 3-2 = 1
\]


\newpage
(δ) \skato Υποθέτω το ερώτημα ζητάει μια βάση για τον χώρο στηλών και μία βάση για τον χώρο γραμμών όταν \textlatin{k=1?}  

\[
\textbf{Έχουμε: }
A =
\begin{bmatrix}
& 1 & 1 & 5 &  \\
& 0 & 0 & -1 & \\
& 0 & 0 & -1 & 
\end{bmatrix}
\Rightarrow
\Vec{c_1} = \begin{bmatrix} 1 \\ 0 \\ 0 \end{bmatrix} =
\skato \Vec{c_2} = \begin{bmatrix} 1 \\ 0 \\ 0 \end{bmatrix},
\skato \vec{c_3} = \begin{bmatrix} 5 \\ -1 \\ -1  \end{bmatrix}
\skata \textbf{Έτσι:} \skata
C(\begin{bmatrix} 1 & 5 \\ 0 & -1 \\ 0 & -1 \end{bmatrix})_{3\times2}
\]

\[
\textbf{Έχουμε: }
A =
\begin{bmatrix}
& 1 & 1 & 5 &  \\
& 0 & 0 & -1 & \\
& 0 & 0 & -1 & 
\end{bmatrix}
\Rightarrow
\Vec{r_1} = \begin{bmatrix} 1 & 1 & 5 \end{bmatrix},
\skato \Vec{r_2} = \begin{bmatrix} 0 & 0 & -1 \end{bmatrix} =
\skato \vec{r_3} = \begin{bmatrix} 0 & 0 & -1 \end{bmatrix}
\skata \textbf{Έτσι:} \skata
C(\begin{bmatrix} 1 & 1 & 5 \\ 0 & 0 & -1 \end{bmatrix})_{2\times3}
\]




% ---------------------------------------------------------------------




\vspace{1cm}
\vspace{1.23456789cm}
\begin{center}
    {{\textbf{·  \underline{2\textsuperscript{η} Άσκηση:}}}}
\end{center}
\vspace{0.54321cm}

\[
\textbf{Έχουμε:} \hspace{0.54321cm} 
A_{3\times4} \skato = \skato
\begin{bmatrix}
& 0 & 1 & 0 & 0 & \\
& 0 & 0 & 1 & 0 & \\
& 1 & 0 & 0 & 0 &
\end{bmatrix}
\skata \textbf{ή} \skata 
A_{3\times4} \skato = \skato
\begin{bmatrix}
& \boxed{1} & 0 & 0 & 0 & \\
& 0 & \boxed{1} & 0 & 0 & \\
& 0 & 0 & \boxed{1} & 0 &
\end{bmatrix}
\skata \Rightarrow \skata
\left\{
\begin{array}{l}
rows(A)=3 \\
columns(A)=4 \\
rank(A) = 3 
\end{array} 
\right.
\]

\[
- \skata dim(C(A))=rank(A)=3 \skata \textbf{και} \skata 
C(\begin{bmatrix}
    0 & 1 & 0 \\ 0 & 0 & 1 \\ 1 & 0 & 0
\end{bmatrix})_{3\times3}
\skata \skata *(c_1, c_2, c_3)
\]

\[
- \skata dim(C(A^t))=rank(A)=3 \skata \textbf{και} \skata 
C(\begin{bmatrix}
& 0 & 1 & 0 & 0 & \\
& 0 & 0 & 1 & 0 & \\
& 1 & 0 & 0 & 0 &
\end{bmatrix})_{3\times4}
= A
\skata \skata *(r_1, r_2, r_3)
\]

\[
- \skata  dim(N(A)) = columns(A) - rank(A) = 4 - 3 = 1
\skata \textbf{και} \skata
Ax=0 \Rightarrow
\begin{bmatrix}
& 0 & 1 & 0 & 0 & \\
& 0 & 0 & 1 & 0 & \\
& 1 & 0 & 0 & 0 &
\end{bmatrix}
\cdot
\begin{bmatrix}
& x_1 & \\
& x_2 & \\
& x_3 & \\
& x_4 &
\end{bmatrix}
=
\begin{bmatrix}
& 0 & \\
& 0 & \\
& 0 & 
\end{bmatrix}
\Rightarrow
\]

\[
\Rightarrow
\begin{bmatrix}
& x_2 & \\
& x_3 & \\
& x_1 & 
\end{bmatrix}
=
\begin{bmatrix}
& 0 & \\
& 0 & \\
& 0 & 
\end{bmatrix}
\Rightarrow
x_1 = x_2 = x_3 = 0
\skata \textbf{και έτσι:} \skata
N(\begin{bmatrix}
& 0 & \\
& 0 & \\
& 0 & \\
& 1 &
\end{bmatrix})_{4\times1}
\skata \skata * (x_4 \textbf{ ελεύθερη})
\]

\[
- \skata  dim(N(A)) = rows(A^t) - rank(A) = 4 - 3 = 1
\skata \textbf{και} \skata 
A^tx'=0 \Rightarrow
\begin{bmatrix}
& 0 & 0 & 1 &  \\
& 1 & 0 & 0 &  \\
& 0 & 1 & 0 &  \\
& 0 & 0 & 0 &
\end{bmatrix}
\cdot
\begin{bmatrix}
& x'_1 & \\
& x'_2 & \\
& x'_3 & 
\end{bmatrix}
=
\begin{bmatrix}
& 0 & \\
& 0 & \\
& 0 & \\
& 0 &
\end{bmatrix}
\Rightarrow
\]

\[
\Rightarrow
\begin{bmatrix}
& x'_3 & \\
& x'_1 & \\
& x'_2 & \\
& 0 &
\end{bmatrix}
=
\begin{bmatrix}
& 0 & \\
& 0 & \\
& 0 & \\
& 0 &
\end{bmatrix}
\Rightarrow
x'_1 = x'_2 = x'_3 = 0
\skata \textbf{και έτσι:} \skata
N(\begin{bmatrix}
& 0 & \\
& 0 & \\
& 0 & \\
& 1 &
\end{bmatrix})_{4\times1}
\]







% ---------------------------------------------------------------------




\vspace{5cm}


\vspace{1.23456789cm}
\begin{center}
    {{\textbf{·  \underline{3\textsuperscript{η} Άσκηση:}}}}
\end{center}
\vspace{0.54321cm}


 \[
\textbf{(α)} \skato \textbf{Πίνακα ώστε: Ο χώρος στηλών να περιέχει τα} \skato (1, 2, -3)^t, (2, -3, 5)^t \skato \textbf{και ο μηδενόχωρος το} \skato (1, 1, 1)^t. 
\]


\[
\textbf{Έστω:} \skata 
c_1=\begin{bmatrix} 1 \\ 2 \\ -3 \end{bmatrix}_{1\times3}, \skato c_2=\begin{bmatrix} 2 \\ -3 \\ 5 \end{bmatrix}_{1\times3}
\skato \Rightarrow \skato
Α = \begin{bmatrix} & c_1 & c_2 & c_3 & \end{bmatrix}
\]

\[
\textbf{Αφού } \skato \vec{x} = \begin{bmatrix} 1\\1\\1 \end{bmatrix}
\epsilon N(A) \skato \Rightarrow Ax = 0 \Rightarrow
\begin{bmatrix} & c_1 & c_2 & c_3 & \end{bmatrix} \cdot
\begin{bmatrix} & 1 & \\& 1 &\\ &1 &\end{bmatrix}=0
\Rightarrow \vec{c_1} + \vec{c_2} + \vec{c_3} = 0
\Rightarrow \skata \vec{c_3} = -(\vec{c_1} + \vec{c_2})
\]

\[
\textbf{Άρα:} \skata 
c_3 = \begin{bmatrix} -3 \\ 1 \\ -2  \end{bmatrix}
\skata \textbf{ και τελικά: } \skata
{
A=
\begin{bmatrix} 
& 1 & 2 & -3 & \\
& 2 & -3 & 1 & \\
& -3 & 5 & -2 &
\end{bmatrix}_{3\times3}
}
\]


\vspace{1cm}


\[
\textbf{(β)} \skato \textbf{Πίνακα ώστε: Ο χώρος γραμμών να περιέχει τα} \skato (1, 2, -3)^t, (2, -3, 5)^t \skato \textbf{και ο μηδενόχωρος το} \skato (1, 1, 1)^t. 
\] 

\[
\textbf{Έστω:} \skata 
r_1=\begin{bmatrix} 1 \\ 2 \\ -3 \end{bmatrix}_{1\times3}, \skato r_2=\begin{bmatrix} 2 \\ -3 \\ 5 \end{bmatrix}_{1\times3}
\skato \Rightarrow \skata
Α = \begin{bmatrix} & r_1^t & \\ & r_2^t & \\ & \cdot \cdot \cdot & \end{bmatrix}_{m\times3}
\]

\[
\textbf{Αφού } \skato \vec{x} = \begin{bmatrix} 1\\1\\1 \end{bmatrix}
\epsilon N(A) \skato \Rightarrow Ax = 0 \Rightarrow
\begin{bmatrix} & r_1^t & \\ & r_2^t & \\ & \cdot \cdot \cdot & \end{bmatrix} \cdot
\begin{bmatrix} & 1 & \\& 1 &\\ & 1 &\end{bmatrix} =
\begin{bmatrix} & 0 & \\& 0 &\\ & \cdot \cdot \cdot &\end{bmatrix}
\Rightarrow 
\begin{bmatrix} & 1 & 2 & -3 & \\ & 2 & -3 & 1 & \\ & & \cdot \cdot \cdot & \end{bmatrix} \cdot
\begin{bmatrix} & 1 & \\& 1 &\\ & 1 &\end{bmatrix} =
\begin{bmatrix} & 0 & \\& 0 &\\ & \cdot \cdot \cdot &\end{bmatrix}
\Rightarrow
\]

\[
\Rightarrow
\begin{bmatrix} & 1 + 2 + (-3) & \\ & 2 + (-3) + 1 & \\ &  \cdot \cdot \cdot & \end{bmatrix} =
\begin{bmatrix} & 0 & \\& 0 &\\ & \cdot \cdot \cdot &\end{bmatrix}
\Rightarrow
\begin{bmatrix} & 0 & \\ & 4 & \\ &  \cdot \cdot \cdot & \end{bmatrix} \neq
\begin{bmatrix} & 0 & \\& 0 &\\ & \cdot \cdot \cdot &\end{bmatrix}
\skata \textbf{αδύνατο, άρα} \skata \vec{x} \notin N(A). 
\]

\[
{\textbf{Τελικά, δεν υπάρχει τέτοιος πίνακας!}}
\]







% ---------------------------------------------------------------------







\newpage
\vspace{1.23456789cm}
\begin{center}
    {{\textbf{·  \underline{4\textsuperscript{η} Άσκηση:}}}}
\end{center}
\vspace{0.1cm}


\[
(αi) \skata \boxed{ S^{\perp} = \mathbb{R}^3 } \hspace{0.25cm}
\textbf{αφού ισχύει:} \skata 
dim(S) + dim(S^{\perp}) = dim(\mathbb{R}) \Rightarrow
dim(S^{\perp}) = dim(\mathbb{R}^3) - dim(S=\{\vec{0}\}) = 3-0 = 3
\]


\[
(αii) \skata 
S=\{(1,1,1)^t\} \skato \textbf{και} \skata
dim(S) = 1 \Rightarrow dim(S^{\perp}) = 3-1=2 \skata
\textbf{(επίπεδο) έτσι:} \skata 
\boxed{S^{\perp} = \{ (x,y,z)^T \in \mathbb{R}^3 \skato | \skato x+y+z=0 \}}
\]


\[
(αiii) \skata 
S=\{(2,0,0)^t, (0,0,3)^t\} \skata \textbf{και} \skata
dim(S) = 2 \Rightarrow dim(S^{\perp}) = 1 \skata
\textbf{(ευθεία) έτσι:} \skata 
\boxed{S^{\perp} = \{ (0,y,0)^T \skato | \skato y \in \mathbb{R} \}}
\]


\[
(βi) \skata 
\boxed{ V^{\perp} = \{ \vec{0} \} } \hspace{0.5cm}
\textbf{αφού ισχύει:} \skata 
dim(V) + dim(V^{\perp}) = dim(\mathbb{R}) \Rightarrow
dim(V^{\perp}) = dim(\mathbb{R}^4) - dim(V) = 4-4 = 0
\]

\[
(βii) \skata 
\boxed{ (V^{\perp})^{\perp} = \mathbb{R}^4 } \hspace{0.5cm}
\textbf{αφού:} \skata
dim((V^{\perp})^{\perp}) = dim(\mathbb{R}^4) - dim((V^{\perp})) = 4-0 = 4 \skata
\]







% ---------------------------------------------------------------------






\vspace{1.23456789cm}
\begin{center}
    {{\textbf{·  \underline{5\textsuperscript{η} Άσκηση:}}}}
\end{center}
\vspace{0.321cm}


Έχουμε πως υποχώρος \textlatin{S} παράγεται από τα διανύσματα:
\skata 
\textlatin{s}1 = \begin{bmatrix} 1 \\ 2 \\ 2 \\ 3 \end{bmatrix}, \skata
\textlatin{s}2 = \begin{bmatrix} 1 \\ 3 \\ 3 \\ 2 \end{bmatrix}

\[
\textbf{Έστω: } \skato
A_{2\times4} = \begin{bmatrix} 1&2&2&3\\1&3&3&2 \end{bmatrix}
\skata \textbf{τότε} \skata Ax=0 \skato \Rightarrow
\begin{bmatrix} 1&2&2&3\\1&3&3&2 \end{bmatrix} \cdot
\begin{bmatrix} x_1\\x_2\\x_3\\x_4 \end{bmatrix}
= \begin{bmatrix} 0\\0 \end{bmatrix}
\skato \overset{Gauss}{\Rightarrow} \skato 
\begin{bmatrix} 1&2&2&3&|&0\\1&3&3&2&|&0 \end{bmatrix} 
\]

\[
=
\begin{bmatrix} \boxed{1}&2&2&3&|&0\\ 0&\boxed{1}&1&-1&|&0 \end{bmatrix} \skato \Rightarrow
\left\{
\begin{array}{l}
1x_1 + 2x_2 + 2x_3 + 3x_4 = 0 \\
0x_1 + 1x_2 + 1x_3 - 1x_4 = 0
\end{array}
\right.
\Rightarrow
\left\{
\begin{array}{l}
1x_1 + 2x_2 + 2x_3 + 3x_4 = 0 \\
x_2 = x_4 -x_3
\end{array}
\right.
\Rightarrow
\left\{
\begin{array}{l}
x_1 = -5x_4 \\
x_2 = x_4 -x_3
\end{array}
\right.
\Rightarrow
\]

\[
\Rightarrow
\left\{
\begin{array}{l}
x_1 = -5x_4 \\
x_2 = -x_3 + x_4 \\
x_3 = x_3 \\
x_4 = x_4
\end{array}
\right.
\skato \Rightarrow \skato
x = x_3 \cdot \begin{bmatrix} 0\\-1\\1\\0 \end{bmatrix} + x_4 \cdot
\begin{bmatrix} -5\\1\\0\\1 \end{bmatrix}
\skata \textbf{ Τελικά, βάση του} \skata S^{\perp}: \skata 
\{ \begin{bmatrix} 0\\-1\\1\\0 \end{bmatrix}, \skato \begin{bmatrix} -5\\1\\0\\1 \end{bmatrix} \} 
\skata \textbf{αφού} \skata N(A)=S^{\perp}.
\]







% ---------------------------------------------------------------------





 
\vspace{1.23456789cm}
\begin{center}
    {{\textbf{·  \underline{6\textsuperscript{η} Άσκηση:}}}}
\end{center}


\[
\textbf{Απόδειξη: } \skata
\textbf{Αν } S \subseteq V \Rightarrow V^\perp \subseteq S^\perp 
\]
\[
\textbf{Έστω } \skato x \in V^\perp \skato \Rightarrow \skato 
x \cdot v = 0 \skato \forall \skato v \in V. \skata \textbf{Όμως, } \skato S \subseteq V \Rightarrow \skato \forall \skato s \in S, v \in V. \skata \textbf{Έτσι, } \skato x \cdot s = 0 \skato \forall \skato s \in S \Rightarrow x \in S^\perp.
\]
\[
\textbf{Τελικά, } \skato \forall \skato x \in V^\perp, x \in S^\perp \skato \Rightarrow \skato V^\perp \subseteq S^\perp. \skata \qed
\]









% ---------------------------------------------------------------------





 
\vspace{3.25cm}
\vspace{1.23456789cm}
\begin{center}
    {{\textbf{·  \underline{7\textsuperscript{η} Άσκηση:}}}}
\end{center}


Έχουμε: \skato \textlatin{3x + y} = 0 \Rightarrow \textlatin{y = -3x}. \textbf{ Άρα, διάνυσμα κατεύθηνσης: } (x,y) = (x, -3x) = x(1, -3) \Rightarrow \vec{u} = \begin{bmatrix}
1 \\ -3 \end{bmatrix}.

 
\[
\textbf{Ξέρουμε πως: } \skato 
P = \frac{(\vec{u}) \cdot (\vec{u}^T)}{(\vec{u}^T) \cdot (\vec{u})}
\skato \textbf{ αν } \skato \vec{u} \in \mathbb{R}^2 \skato \textbf{ ο πίνακας προβολής.}
\]

\[
\textbf{Άρα: } \skato 
P=\frac{\begin{bmatrix}1\\-3\end{bmatrix} \cdot \begin{bmatrix}1 & -3\end{bmatrix}}{\begin{bmatrix}1&-3\end{bmatrix} \cdot \begin{bmatrix}1\\-3\end{bmatrix}}
\skato = \skato 
\frac{\begin{bmatrix}1&-3\\-3&9\end{bmatrix}}{\begin{bmatrix}(1)(1)+(-3)(-3)\end{bmatrix}}
\skato = \skato
\frac{1}{10} \cdot \begin{bmatrix}1&-3\\-3&9\end{bmatrix}
\skato = \skato
{
\begin{bmatrix}1/10 & -3/10 \\ -3/10 & 9/10\end{bmatrix}
\skato = \skato P_{2\times2}
}
\]






% ---------------------------------------------------------------------







 \vspace{0.1cm}
\vspace{1.23456789cm}
\begin{center}
    {{\textbf{·  \underline{8\textsuperscript{η} Άσκηση:}}}}
\end{center}
\vspace{0.01cm}


\[
\textbf{Έχουμε }: \skata α=\begin{bmatrix} 1\\1\\0 \end{bmatrix}, \skato
β=\begin{bmatrix} 1\\0\\1 \end{bmatrix}, \skato
c=\begin{bmatrix} 0\\1\\1 \end{bmatrix}.
\skata \textbf{Ψάχνουμε }: \skato 
\vec{u_1}=\frac{a}{||a||}, \skato 
\vec{u_2}=\frac{b'}{||b'||}, \skato 
\vec{u_3}=\frac{c''}{||c''||}
\]

\[
* \skata 
||a|| = \sqrt{1^2+1^2+0^2}= \sqrt{2} \skato \Rightarrow \skato
\vec{u_1} = \frac{\begin{bmatrix} 1\\1\\0 \end{bmatrix}}{\sqrt{2}}=
\frac{1}{\sqrt{2}} \cdot \begin{bmatrix} 1\\1\\0 \end{bmatrix} =
{\begin{bmatrix} 1/\sqrt{2}\\1/\sqrt{2}\\0 \end{bmatrix}=
\underline{ \vec{u_1}_{(3\times1)}} }
\]

\[
* \skata
b' = b - proj_{u_1}(b) = b - \frac{b \cdot u_1}{||u_1||^2} \cdot u_1 = \begin{bmatrix} 1\\0\\1 \end{bmatrix} -
\frac{\begin{bmatrix} 1\\0\\1 \end{bmatrix} \cdot \begin{bmatrix} 1/\sqrt{2}\\1/\sqrt{2}\\0 \end{bmatrix}}{\sqrt{(1/\sqrt{2})^2+(1/\sqrt{2})^2+0^2}^2} \cdot \begin{bmatrix} 1/\sqrt{2}\\1/\sqrt{2}\\0 \end{bmatrix}
= \begin{bmatrix} 1\\0\\1 \end{bmatrix} - \frac{1}{\sqrt{2}} \cdot \begin{bmatrix} 1/\sqrt{2}\\1/\sqrt{2}\\0 \end{bmatrix}
= \begin{bmatrix} 1\\0\\1 \end{bmatrix} -
\begin{bmatrix} 1/2\\1/2\\0 \end{bmatrix}
=
\]
\[
\Rightarrow b' = \begin{bmatrix} 1/2\\-1/2\\1 \end{bmatrix}
\skato \textbf{kai} \skato 
u_2=\frac{b'}{||b'||}=
\frac{\begin{bmatrix} 1/2\\-1/2\\1 \end{bmatrix}}
{\sqrt{(1/2)^2+(-1/2)^2+(1)^2}} = \frac{1}{\frac{\sqrt{6}}{2}} \cdot \begin{bmatrix} 1/2\\-1/2\\1 \end{bmatrix} =
\frac{2}{\sqrt{6}} \cdot \begin{bmatrix} 1/2\\-1/2\\1 \end{bmatrix}
= { \begin{bmatrix} 1/\sqrt{6}\\-1/\sqrt{6}\\2/\sqrt{6} \end{bmatrix} = 
\underline{ u_2_{(3\times1)}} }
\]

\[
* \skata
c'' = c - proj_{u_1}(c) - proj_{u_2}(c) =
c - \frac{c \cdot u_1}{||u_1||^2} \cdot u_1 - \frac{c \cdot u_2}{||u_2||^2} \cdot u_2 =
\begin{bmatrix} 0\\1\\1 \end{bmatrix} - \frac{\begin{bmatrix} 0\\1\\1 \end{bmatrix} \cdot \begin{bmatrix} 1/\sqrt{2}\\1/\sqrt{2}\\0 \end{bmatrix}}{1} \cdot \begin{bmatrix} 1/\sqrt{2}\\1/\sqrt{2}\\0 \end{bmatrix}
-  
\]
\[
- \frac{\begin{bmatrix} 0\\1\\1 \end{bmatrix} \cdot \begin{bmatrix} 1/\sqrt{6}\\-1/\sqrt{6}\\2/\sqrt{6} \end{bmatrix}}{\sqrt{(1/\sqrt{6})^2+(-1/\sqrt{6)^2}+(2/\sqrt{6})^2}} \cdot \begin{bmatrix} 1/\sqrt{6}\\-1/\sqrt{6}\\2/\sqrt{6} \end{bmatrix}
= 
\begin{bmatrix} 0\\1\\1 \end{bmatrix}
- \frac{1}{\sqrt{2}} \cdot \begin{bmatrix} 1/\sqrt{2}\\1/\sqrt{2}\\0 \end{bmatrix} - \frac{\frac{1}{\sqrt{6}}}{1} \cdot \begin{bmatrix} 1/\sqrt{6}\\-1/\sqrt{6}\\2/\sqrt{6} \end{bmatrix}
=
\begin{bmatrix} 0\\1\\1 \end{bmatrix} - 
\begin{bmatrix} 1/2\\1/2\\0 \end{bmatrix} -
\begin{bmatrix} 1/6\\-1/6\\1/3 \end{bmatrix}
=
\]
\[
= \begin{bmatrix} 0-1/2-1/6\\1-1/2-(-1/6)\\1-0-1/3 \end{bmatrix}
= {\begin{bmatrix} -2/3\\2/3\\2/3 \end{bmatrix} = c'' }
\skato \textbf{kai} \skato
u_3 = \frac{\begin{bmatrix} -2/3\\2/3\\2/3 \end{bmatrix}}{\sqrt{(-2/3)^2+(2/3)^2+(2/3)^2}} =
\frac{\sqrt{3}}{2} \cdot \begin{bmatrix} -2/3\\2/3\\2/3 \end{bmatrix} =
{\begin{bmatrix} -\sqrt{3}/3\\ \sqrt{3}/3\\ \sqrt{3}/3 \end{bmatrix} = 
\underline{ u_3_{(3\times1)}} }
\]







% ---------------------------------------------------------------------





 
\vspace{0cm}
\vspace{1.23456789cm}
\begin{center}
    {{\textbf{·  \underline{9\textsuperscript{η} Άσκηση:}}}}
\end{center}




\[
\textbf{Έχουμε:} \skata 
Α=
\begin{bmatrix}
    & 1 & 1 & 0 & \\
    & 1 & 0 & 1 & \\
    & 0 & 1 & 1 &
\end{bmatrix}
\skata 
\textbf{και}
\skata
A=3\times3 \skato \longrightarrow \skato Q=3\times3, \skato R=3\times3.
\]

\[
\textbf{Από προηγούμενη άσκηση:} \skata
c_1^A= \begin{bmatrix} 1\\1\\0 \end{bmatrix}
\overset{Gram}{\underset{Schmidt}{\Rightarrow}}
\begin{bmatrix} \frac{1}{\sqrt{2}}\\ \frac{1}{\sqrt{2}}\\0 \end{bmatrix}, \skata 
c_2^A= \begin{bmatrix} 1\\0\\1 \end{bmatrix}
\overset{Gram}{\underset{Schmidt}{\Rightarrow}}
\begin{bmatrix} \frac{1}{\sqrt{6}}\\ -\frac{1}{\sqrt{6}}\\ \frac{2}{\sqrt{6}} \end{bmatrix}, \skata
c_3^A = \begin{bmatrix} 0\\1\\1 \end{bmatrix}
\overset{Gram}{\underset{Schmidt}{\Rightarrow}}
\begin{bmatrix} -\frac{\sqrt{3}}{3}\\ \frac{\sqrt{3}}{3}\\ \frac{\sqrt{3}}{3} \end{bmatrix}
\]

\[
\textbf{Έτσι:} \skata
Q_{3\times3} \skato = \skato
\begin{bmatrix} 
\frac{1}{\sqrt{2}} & \frac{1}{\sqrt{6}} & -\frac{\sqrt{3}}{3} \\
\frac{1}{\sqrt{2}} & -\frac{1}{\sqrt{6}} & \frac{\sqrt{3}}{3} \\
0 & \frac{2}{\sqrt{6}} & \frac{\sqrt{3}}{3}
\end{bmatrix}
\]

\[
\textbf{Και, αφού:} \skata
A = Q \cdot R \Rightarrow R = Q^T \cdot A = 
\begin{bmatrix} 
\frac{1}{\sqrt{2}} & \frac{1}{\sqrt{2}} & 0 \\
\frac{1}{\sqrt{6}} & -\frac{1}{\sqrt{6}} & \frac{2}{\sqrt{6}} \\
-\frac{\sqrt{3}}{3} & \frac{\sqrt{3}}{3} & \frac{\sqrt{3}}{3}
\end{bmatrix}
\cdot
\begin{bmatrix}
    & 1 & 1 & 0 & \\
    & 1 & 0 & 1 & \\
    & 0 & 1 & 1 &
\end{bmatrix}
=
\]
\[
=
\begin{bmatrix}
    \frac{1}{\sqrt{2}} \cdot1 + \frac{1}{\sqrt{2}} \cdot1 + 0\cdot0 & \frac{1}{\sqrt{2}} \cdot1 + \frac{1}{\sqrt{2}} \cdot0 + 0 \cdot 1 & \frac{1}{\sqrt{2}} \cdot0 + \frac{1}{\sqrt{2}} \cdot1 + 0 \cdot 1 \\
    \frac{1}{\sqrt{6}} \cdot1 + (-\frac{1}{\sqrt{6}}) \cdot1 + \frac{2}{\sqrt{6}} \cdot0 & \frac{1}{\sqrt{6}} \cdot1 + (-\frac{1}{\sqrt{6}} \cdot0) + \frac{2}{\sqrt{6}} \cdot1 & \frac{1}{\sqrt{6}} \cdot0 + (-\frac{1}{\sqrt{6}}) \cdot1 + \frac{2}{\sqrt{6}} \cdot1 \\
    (-\frac{\sqrt{3}}{3} \cdot1 + \frac{\sqrt{3}}{3} \cdot1 + \frac{\sqrt{3}}{3} \cdot0 & (-\frac{\sqrt{3}}{3}) \cdot1 + \frac{\sqrt{3}}{3} \cdot0 + \frac{\sqrt{3}}{3} \cdot1 & (-\frac{\sqrt{3}}{3}) \cdot0 + \frac{\sqrt{3}}{3} \cdot1 + \frac{\sqrt{3}}{3} \cdot1
    \end{bmatrix}
\skato
=
\skato
\begin{bmatrix}
    & \sqrt{2} & \frac{1}{\sqrt{2}} & \frac{1}{\sqrt{2}} & \\
    & 0 & \frac{\sqrt{6}}{2} & \frac{\sqrt{6}}{6} & \\
    & 0 & 0 & \frac{2\sqrt{3}}{3} &
\end{bmatrix}
= R
\]

\[
\textbf{Τελικά:} \skata
{ \skato
A \skato = \skato
\begin{bmatrix} 
\frac{1}{\sqrt{2}} & \frac{1}{\sqrt{6}} & -\frac{\sqrt{3}}{3} \\
\frac{1}{\sqrt{2}} & -\frac{1}{\sqrt{6}} & \frac{\sqrt{3}}{3} \\
0 & \frac{2}{\sqrt{6}} & \frac{\sqrt{3}}{3}
\end{bmatrix}^Q
\skato \cdot \skato
\begin{bmatrix}
    & \sqrt{2} & \frac{1}{\sqrt{2}} & \frac{1}{\sqrt{2}} & \\
    & 0 & \frac{\sqrt{6}}{2} & \frac{\sqrt{6}}{6} & \\
    & 0 & 0 & \frac{2\sqrt{3}}{3} &
\end{bmatrix}^R
\skato }
\]







% ---------------------------------------------------------------------






 \newpage
\vspace{1.23456789cm}
\begin{center}
    {{\textbf{·  \underline{Άσκηση Προγραμματισμού:}}}}
\end{center}
\vspace{0.54321cm}



\setlength{\fboxrule}{1pt}
\begin{figure}[h]
  \hspace{-0.3cm}
  \centering
  \fbox{\includegraphics[width=0.75\textwidth]{images.png}}
\end{figure}


\vspace{0.75cm}

* Γιατί λείπουν ακόμα μερικά εικονοστοιχεία στον χώρο στηλών ?
Παρόλο που ο υπολογισμός του χώρου στηλών αναπληρώνει αρκετά από τα κενά pixels,
εξακολουθούν να υπάρχουν ελλείψεις, διότι το column space περιλαμβάνει μόνο τα
διανύσματα (στήλες) που είναι γραμμικά ανεξάρτητα. Αν ορισμένες πληροφορίες των
αρχικών pixels ήταν εξαρχής χαμένες ή ασυμπλήρωτες, ο χώρος στηλών δεν μπορεί να
τις ανακατασκευάσει τέλεια, καθώς δεν υπάρχει αρκετή πληροφορία για να αναπαραχθεί
όλο το σχήμα της εικόνας.

\vspace{0.5cm}

* Ποια είναι η διάσταση του χώρου στηλών για κάθε εικόνα ?
Για την πρώτη εικόνα \textlatin{(one-a.png)}, η διάσταση του χώρου στηλών είναι 8.
Για τη δεύτερη εικόνα \textlatin{(three.png)}, η διάσταση του χώρου στηλών είναι 9.
Για την τρίτη εικόνα \textlatin{(one-b.png)}, η διάσταση του χώρου στηλών είναι 4.
Για την τέταρτη και τελευταία εικόνα \textlatin{(two.png)}, η διάσταση του χώρου στηλών είναι 13.

\vspace{0.5cm}

* Εάν υπολογίζατε τους χώρους γραμμών, τι θα συνέβαινε ? Θα διορθώνονταν τα κενά ?
Εάν υπολογίζαμε τον χώρο γραμμών (row space) αντί για τον χώρο στηλών, θα
αναπληρώναμε την πληροφορία σε οριζόντιο επίπεδο, δηλαδή κατά μήκος των γραμμών
της εικόνας. Είναι πιθανό ότι ορισμένα κενά pixels θα διορθώνονταν καλύτερα σε
γραμμές που έχουν παρόμοια μοτίβα. Ωστόσο, γενικά το αποτέλεσμα δεν θα ήταν τόσο
καλό όσο με τον χώρο στηλών, διότι η δομή των αριθμών (όπως το σχήμα του ψηφίου)
τείνει να διατηρείται πιο πιστά στις στήλες και όχι τόσο στις γραμμές.




 

% ---------------------------------------------------------------------------------------


 \vspace{10cm}

\begin{center}
    \vspace*{\fill}
    \setlength{\fboxsep}{10pt}
    \boxed{\textbf{\LARGE{ΤΕΛΟΣ}}}
    \vspace*{\fill}
\end{center}

\end{document}


% ---------------------------------------------------------------------------------------


 
% The end ...